\chapter{Tabular Methods}

\begin{comment}
    Short description of what tabular methods are.
     - They keep the value function in a table
     - Discreet states and actions

    Explain k-armed bandit agent from previous chapter.

    Upper confidence bound

    Dynamic programming

    Monte carlo methods

    Temporal difference
     - Sarsa
     - Expected Sarsa
     - Q-learning
     - Double Q-learning

    n-step bootstrapping

    Planning and learning (maybe not)
     - models and planning
     - dyna

    Policy gradient
    
    Working example(s)
     - Discreet state and action space
     - model
       • quadcopter navigation
       • pole
       • some other thing I haven't made yet

\end{comment}

The action-value function estimate $Q_t(s, a)$ holds the current estimate of the real action-value function $q_*(a)$. If the \gls{rl}-agent is to achieve it's goal, $Q_t(s, a)$ needs to be as close to $q_*(a)$ as possible. The actual function used to represent $Q_t(s, a)$ can be different depending on the application, and the methods used to estimate $Q_t(s, a)$ can be different depending on the underling function and application. 

This chapter will focus on methods where $q_*(a)$ is estimated using tables of discreet values. These methods are useful with small discreet state- and action-spaces and can be used to explain different concepts.
